\subsection*{導入：設計とは何か}
\addcontentsline{toc}{subsection}{導入：設計とは何か}

ソフトウェア設計という語は、近いが異なる意味で使われる。第一に、設計は分野である。これは、科学的原理、技術情報、想像力を使って、目的を満たすソフトウェアを経済的かつ効率的に定義する分野である。第二に、設計はその分野の中で行われるプロセスである。第三に、設計はプロセスの結果、つまり設計記述やモデルである。第四に、設計はライフサイクル上の段階でもある。

\term{ソフトウェア設計記述}{Software Design Description, SDD} は、設計情報を伝えるための成果物である。これは、実装、計画、分析、意思決定を助ける「青写真」に相当する。設計記述には、ソフトウェアを構成要素へ分けること、構成要素をどのように配置するか、構成要素同士と外部世界とのインタフェースをどのように定義するかが含まれる。

ソフトウェア設計は、要求を分析し、外部から見える性質と内部構造を定める活動である。設計は、一般に三段階で考えると理解しやすい。第一はアーキテクチャ設計であり、システム全体の根本的な構造と原則を扱う。第二は高水準設計であり、システムと主要構成要素が外部とどう関わるかを扱う。第三は詳細設計であり、各構成要素の内部を実装できる水準まで具体化する。

\begin{pointbox}{重要}
\term{設計}{Design}は「実装前にきれいな図を描くこと」ではない。設計は、要求、制約、品質目標、技術選択、保守の見通しを踏まえ、作れる解決策を定義する活動である。
\end{pointbox}
